% 08_twocolumn_siggraph.tex
% Two-column Siggraph format (ACM SIGGRAPH proceedings style)
% Overleaf compatible: YES (acmart class)

% ACM SIGGRAPH template
% Use: acmart class with sigconf option
% For actual submission, use the official ACM template from:
% https://www.acm.org/publications/proceedings-template

\documentclass[sigconf,review,anonymous]{acmart}
% Options:
% sigconf    - SIGGRAPH conference format (two-column)
% review     - adds line numbers for review
% anonymous  - anonymizes authors for double-blind review
% Non-anonymous: remove 'anonymous' and add author info below

\usepackage[utf8]{inputenc}
\usepackage[T1]{fontenc}
\usepackage{lmodern}
\usepackage[english]{babel}
\usepackage{hyperref}
\usepackage{cleveref}
\usepackage{graphicx}
\usepackage{subcaption}
\usepackage{booktabs}
\usepackage{amsmath}
\usepackage{amssymb}
\usepackage{mathtools}
\usepackage{physics}
\usepackage{algorithm2e}
\usepackage{algpseudocode}
\usepackage{listings}
\usepackage{xcolor}

% ACM metadata (fill in for final submission)
\acmConference[SIGGRAPH '26]{ACM SIGGRAPH 2026 Conference Papers}{August 2026}{Denver, CO, USA}
\acmYear{2026}
\acmISBN{978-1-4503-XXXX-X/26/08}
\acmDOI{10.1145/XXXXXXX.XXXXXXX}

% Copyright (choose one)
%\setcopyright{acmcopyright}
%\setcopyright{acmlicensed}
%\setcopyright{rightsretained}
%\setcopyright{usgov}
%\setcopyright{usgovmixed}
%\setcopyright{cusgov}
%\setcopyright{cusgovmixed}

% For review mode - line numbers
\linenumbers

% Algorithm2e setup
\SetAlgoLined
\SetKwInOut{Input}{Input}
\SetKwInOut{Output}{Output}

\hypersetup{
    colorlinks=true,
    linkcolor=blue,
    citecolor=red,
    urlcolor=cyan
}

% Paper title
\title{Two-Column SIGGRAPH Paper Template}

% Authors (anonymous for review)
%\author{Anonymous Author(s)}
%\affiliation{\institution{Anonymous Institution}}

% For non-anonymous version:
%\author{First Author}
%\authornote{Equal contribution}
%\affiliation{\institution{University A}}
%\email{first@university-a.edu}
%
%\author{Second Author}
%\affiliation{\institution{Company B}}
%\email{second@company-b.com}

\begin{document}

\maketitle

\begin{abstract}
This template demonstrates the two-column ACM SIGGRAPH format
using the \texttt{acmart} class with \texttt{sigconf} option.
It includes all standard sections, figures, tables, algorithms,
and references formatted for SIGGRAPH proceedings.
\end{abstract}

\begin{CCSXML}
<ccs2012>
<concept>
<concept_id>10010147.10010371.10010352.10010379</concept_id>
<concept_desc>Computing methodologies~Rendering</concept_desc>
<concept_significance>500</concept_significance>
</concept>
<concept>
<concept_id>10010147.10010371.10010389</concept_id>
<concept_desc>Computing methodologies~Ray tracing</concept_desc>
<concept_significance>400</concept_significance>
</concept>
</ccs2012>
\end{CCSXML}

\ccsdesc[500]{Computing methodologies~Rendering}
\ccsdesc[400]{Computing methodologies~Ray tracing}

\keywords{rendering, ray tracing, neural networks, real-time graphics}

\begin{teaserfigure}
\centering
\fbox{\parbox{0.9\linewidth}{\centering [Teaser Figure: Main result image spanning both columns]}}
\caption{Teaser figure showing main result. This spans both columns
using the \texttt{teaserfigure} environment.}\label{fig:teaser}
\end{teaserfigure}

\section{Introduction}
\label{sec:intro}

SIGGRAPH papers use the two-column \texttt{acmart} format.
This template provides a starting point for conference submissions.

Key formatting requirements:
\begin{itemize}
    \item Two-column layout (\texttt{sigconf} option)
    \item 8-10 pages for full papers (excluding references)
    \item Anonymous submission for double-blind review
    \item ACM CCS categories and keywords
    \item Teaser figure spanning both columns
\end{itemize}

\section{Related Work}
\label{sec:related}

Cite prior work using \texttt{biblatex} or \texttt{BibTeX} with ACM style.
Example citations: \cite{knuth1984texbook}, \cite{whitted1980improved},
\cite{akanine2018realtime}.

\section{Method}
\label{sec:method}

\subsection{Overview}
\label{sec:overview}

Our method consists of three stages:
\begin{enumerate}
    \item Stage 1: Description
    \item Stage 2: Description
    \item Stage 3: Description
\end{enumerate}

\subsection{Technical Details}
\label{sec:technical}

Mathematical formulation:

\begin{equation}\label{eq:method}
L_o(\vb{x}, \omega_o) = \int_{\Omega} f_r(\vb{x}, \omega_i, \omega_o) L_i(\vb{x}, \omega_i) \abs{\vb{n} \vdot \omega_i} \d\omega_i
\end{equation}

Algorithm pseudocode:

\begin{algorithm}[htbp]
\caption{Our Algorithm}\label{alg:ours}
\Input{Input data}
\Output{Result}
\ForEach{element}{
    process\;
}
\Return result
\end{algorithm}

\subsection{Implementation}
\label{sec:impl}

GPU implementation details. Use \texttt{listings} for code:

\begin{lstlisting}[language=C, caption=CUDA kernel, label=lst:cuda]
__global__ void kernel(float* data, int n) {
    int idx = blockIdx.x * blockDim.x + threadIdx.x;
    if (idx < n) data[idx] *= 2.0f;
}
\end{lstlisting}

\section{Results}
\label{sec:results}

\subsection{Quantitative Evaluation}
\label{sec:quantitative}

\begin{table}[htbp]
\caption{Quantitative comparison}\label{tab:quant}
\begin{tabular}{lccc}
\toprule
Method & PSNR $\uparrow$ & SSIM $\uparrow$ & Time (ms) \\
\midrule
Baseline & 28.5 & 0.92 & 45 \\
Ours & \textbf{32.1} & \textbf{0.96} & 52 \\
\bottomrule
\end{tabular}
\end{table}

\subsection{Qualitative Results}
\label{sec:qualitative}

\begin{figure}[htbp]
\centering
\begin{subfigure}[b]{0.48\linewidth}
\centering
\fbox{\parbox{4cm}{\centering [Input]}}
\caption{Input}\label{subfig:input}
\end{subfigure}
\hfill
\begin{subfigure}[b]{0.48\linewidth}
\centering
\fbox{\parbox{4cm}{\centering [Ours]}}
\caption{Ours}\label{subfig:ours}
\end{subfigure}
\caption{Qualitative comparison}\label{fig:qual}
\end{figure}

\section{Discussion}
\label{sec:discussion}

Limitations and future work.

\section{Conclusion}
\label{sec:conclusion}

Summary of contributions.

\section{Acknowledgments}
\label{sec:ack}

We thank reviewers. This work was supported by Grant XXX.

\bibliographystyle{ACM-Reference-Format}
\bibliography{references}

\end{document}